\paragraph{末位二态分层、逆阈值几何与 Bernoulli 塔的 Mellin 编码}
在稳定层均匀经验谱的两点极限、escort 温度族的末位 Bernoulli 塌缩、中心临界容量骨架的双折点闭式以及实参数矩谱与 Bernoulli 模态系数之间的桥接已经建立之后，还可把阈值/采样读出、路径/相位重加权与 Stirling--Bernoulli 层级进一步压缩为同一条末位二态主线。下列结果分别刻画最终严格分层、分位二原子塌缩、归一化逆阈值几何、escort 熵谱的完整常数项，以及容量曲线对 Bernoulli 塔几何权重的 Mellin 编码。

\begin{theorem}[末位相位的最终严格分层]\label{thm:xi-time-part70ab-terminal-bit-eventual-strict-stratification}
\leanverified{paper\_xi\_time\_part70ab\_terminal\_bit\_eventual\_strict\_stratification}
沿用记号
\[
d_m(w):=d_m^{\mathrm{bin}}(w),\qquad
A_i^{(m)}:=\{w\in X_m:\ w_m=i\}\quad(i=0,1),
\qquad
A_m:=\left(\frac{2}{\varphi}\right)^m.
\]
则存在整数 \(m_0\)，使得对一切 \(m\ge m_0\)，都有严格分离
\[
\min_{w\in A_0^{(m)}} d_m(w)
>
\max_{v\in A_1^{(m)}} d_m(v).
\]
因此在充分高分辨率下，纤维多重度的全序被末位 \(w_m\) 强制划分为两层。
\end{theorem}
\begin{proof}
由定理 \ref{thm:fold-bin-two-state-asymptotic}，存在绝对常数 \(C>0\)，使对充分大的 \(m\) 与全部 \(x\in X_m\) 都有
\[
\bigl|d_m(x)-\varphi^{-x_m}A_m\bigr|\le C.
\]
若 \(w\in A_0^{(m)}\) 且 \(v\in A_1^{(m)}\)，则
\[
d_m(w)-d_m(v)
\ge
A_m-C-\bigl(\varphi^{-1}A_m+C\bigr)
=
\varphi^{-2}A_m-2C.
\]
由于 \(A_m\to\infty\)，右端对充分大的 \(m\) 为正，结论成立。
\end{proof}

\begin{corollary}[极值纤维的末位定位与最小尺度失配率]\label{cor:xi-time-part70ab-extremal-fibers-terminal-bit-localization}
\leanverified{paper\_xi\_time\_part70ab\_extremal\_fibers\_terminal\_bit\_localization}
定义
\[
d_{\min}(m):=\min_{w\in X_m}d_m(w),
\qquad
M_m:=\max_{w\in X_m}d_m(w).
\]
则对充分大的 \(m\)，有
\[
\arg\min_{w\in X_m} d_m(w)\subseteq A_1^{(m)},
\qquad
\arg\max_{w\in X_m} d_m(w)\subseteq A_0^{(m)}.
\]
并且
\[
d_{\min}(m)=\varphi^{-1}A_m+O(1),
\qquad
M_m=A_m+O(1),
\qquad
\frac{M_m}{d_{\min}(m)}\longrightarrow \varphi.
\]
此外，沿偶数分辨率 \(m\to\infty\) 有
\[
\frac{d_{\min}(m)}{F_{m/2}}\longrightarrow 0.
\]
\end{corollary}
\begin{proof}
由定理 \ref{thm:xi-time-part70ab-terminal-bit-eventual-strict-stratification}，对充分大的 \(m\)，最小值只能取自 \(A_1^{(m)}\)，最大值只能取自 \(A_0^{(m)}\)。另一方面，由定理 \ref{thm:fold-bin-two-state-asymptotic} 的一致渐近，
\[
d_m(w)=\varphi^{-1}A_m+O(1)\quad(w\in A_1^{(m)}),
\qquad
d_m(w)=A_m+O(1)\quad(w\in A_0^{(m)}),
\]
且误差在各块上统一，故
\[
d_{\min}(m)=\varphi^{-1}A_m+O(1),
\qquad
M_m=A_m+O(1).
\]
于是 \(M_m/d_{\min}(m)\to \varphi\)。

再沿偶数分辨率比较 \(F_{m/2}\)。由 Binet 公式，
\[
F_{m/2}\sim \frac{\varphi^{m/2}}{\sqrt5},
\qquad
d_{\min}(m)\sim \varphi^{-1}\left(\frac{2}{\varphi}\right)^m,
\]
从而
\[
\frac{d_{\min}(m)}{F_{m/2}}
\sim
\sqrt5\,\varphi^{-1}
\left(\frac{2}{\varphi^{3/2}}\right)^m.
\]
由于 \(\varphi^3=2+\sqrt5>4\)，故 \(\varphi^{3/2}>2\)，于是上式趋于 \(0\)。
\end{proof}

\begin{theorem}[中心投影的分位二原子塌缩]\label{thm:xi-time-part70ab-uniform-quantile-two-atom-collapse}
\leanverified{paper\_xi\_time\_part70ab\_uniform\_quantile\_two\_atom\_collapse}
设 \(U_m\) 服从 \(X_m\) 上的均匀分布，并定义尺度化随机变量
\[
Y_m:=\left(\frac{\varphi}{2}\right)^m d_m(U_m).
\]
再定义左连续分位函数
\[
Q_m(\alpha):=\inf\{t\ge 0:\ \mathbf P(Y_m\le t)\ge \alpha\},
\qquad
0<\alpha<1.
\]
则对每个
\[
\alpha\in(0,1)\setminus\{\varphi^{-2}\},
\]
都有极限
\[
Q_m(\alpha)\longrightarrow Q(\alpha),
\qquad
Q(\alpha)=
\begin{cases}
\varphi^{-1}, & 0<\alpha<\varphi^{-2},\\[2mm]
1, & \varphi^{-2}<\alpha<1.
\end{cases}
\]
并且在临界分位 \(\alpha=\varphi^{-2}\) 处出现奇偶振荡：
\[
Q_{2n}(\varphi^{-2})\longrightarrow 1,
\qquad
Q_{2n+1}(\varphi^{-2})\longrightarrow \varphi^{-1}.
\]
\end{theorem}
\begin{proof}
取
\[
0<\varepsilon<\frac{1-\varphi^{-1}}{3}.
\]
由定理 \ref{thm:fold-bin-two-state-asymptotic}，对充分大的 \(m\) 与全部 \(w\in X_m\) 都有
\[
\left|\left(\frac{\varphi}{2}\right)^m d_m(w)-\varphi^{-w_m}\right|
\le
C\left(\frac{\varphi}{2}\right)^m<\varepsilon.
\]
因此当 \(m\) 充分大时，
\[
Y_m\in[\varphi^{-1}-\varepsilon,\varphi^{-1}+\varepsilon]
\quad\text{在 }A_1^{(m)}\text{ 上},
\]
\[
Y_m\in[1-\varepsilon,1+\varepsilon]
\quad\text{在 }A_0^{(m)}\text{ 上}.
\]
故分位值只由块质量
\[
\mathbf P\bigl(Y_m\le \varphi^{-1}+\varepsilon\bigr)
=
\frac{|A_1^{(m)}|}{|X_m|}
=
\frac{F_m}{F_{m+2}}
\]
决定。若 \(\alpha<\varphi^{-2}\) 或 \(\alpha>\varphi^{-2}\)，由
\[
\frac{F_m}{F_{m+2}}\longrightarrow \varphi^{-2}
\]
知对充分大的 \(m\)，比较关系 \(\alpha\lessgtr F_m/F_{m+2}\) 不再改变，于是 \(Q_m(\alpha)\) 分别落在左原子簇与右原子簇中，令 \(\varepsilon\to 0\) 即得所示极限。

最后，由 Binet 公式可直接计算
\[
\frac{F_m}{F_{m+2}}-\varphi^{-2}
=
\frac{(-1)^{m+1}(1-\varphi^{-4})\varphi^{-m}}{\sqrt5\,F_{m+2}},
\]
其符号恰为 \((-1)^{m+1}\)。因此当 \(m\) 为奇数时，
\[
\frac{F_m}{F_{m+2}}>\varphi^{-2},
\]
故 \(Q_m(\varphi^{-2})\) 落在左原子簇；当 \(m\) 为偶数时，
\[
\frac{F_m}{F_{m+2}}<\varphi^{-2},
\]
故 \(Q_m(\varphi^{-2})\) 落在右原子簇。令 \(\varepsilon\to 0\) 即得两条子序列极限。
\end{proof}

\begin{theorem}[escort 温度族的分位相图]\label{thm:xi-time-part70ab-escort-quantile-phase-diagram}
对任意 \(q\ge 0\)，令
\[
W_m^{(q)}\sim \pi_{m,q},
\qquad
Y_m^{(q)}:=\left(\frac{\varphi}{2}\right)^m d_m\bigl(W_m^{(q)}\bigr),
\qquad
\theta(q):=\frac{1}{1+\varphi^{q+1}}.
\]
并定义左连续分位函数
\[
Q_{m,q}(\alpha):=\inf\{t\ge 0:\ \mathbf P(Y_m^{(q)}\le t)\ge \alpha\},
\qquad
0<\alpha<1.
\]
则对每个固定 \(q\ge 0\) 与每个
\[
\alpha\in(0,1)\setminus\{\theta(q)\},
\]
都有极限
\[
Q_{m,q}(\alpha)\longrightarrow Q_q(\alpha),
\qquad
Q_q(\alpha)=
\begin{cases}
\varphi^{-1}, & 0<\alpha<\theta(q),\\[2mm]
1, & \theta(q)<\alpha<1.
\end{cases}
\]
\end{theorem}
\begin{proof}
同样取
\[
0<\varepsilon<\frac{1-\varphi^{-1}}{3}.
\]
由定理 \ref{thm:fold-bin-two-state-asymptotic}，对充分大的 \(m\) 与全部 \(w\in X_m\) 都有
\[
\left|\left(\frac{\varphi}{2}\right)^m d_m(w)-\varphi^{-w_m}\right|<\varepsilon.
\]
因此当 \(m\) 充分大时，
\[
Y_m^{(q)}\in[\varphi^{-1}-\varepsilon,\varphi^{-1}+\varepsilon]
\quad\text{在事件 }W_m^{(q)}(m)=1\text{ 上},
\]
\[
Y_m^{(q)}\in[1-\varepsilon,1+\varepsilon]
\quad\text{在事件 }W_m^{(q)}(m)=0\text{ 上}.
\]
另一方面，由命题 \ref{prop:fold-bin-escort-lastbit}，
\[
\mathbf P\bigl(W_m^{(q)}(m)=1\bigr)=\theta(q)+O\!\left(\left(\frac{\varphi}{2}\right)^m\right).
\]
于是当 \(\alpha<\theta(q)\) 时，充分大的 \(m\) 上有 \(\alpha<\mathbf P(W_m^{(q)}(m)=1)\)，故 \(Q_{m,q}(\alpha)\) 落在左原子簇；当 \(\alpha>\theta(q)\) 时，则 \(Q_{m,q}(\alpha)\) 落在右原子簇。令 \(\varepsilon\to 0\) 即得结论。
\end{proof}

\begin{theorem}[中心容量逆问题的闭式解]\label{thm:xi-time-part70ab-center-capacity-inverse-threshold-closed}
\leanverified{paper\_xi\_time\_part70ab\_center\_capacity\_inverse\_threshold\_closed}
定义归一化连续容量
\[
\Psi_m(\tau):=
2^{-m}\,\mathcal C_m^{\mathrm{bin,cont}}(\tau A_m),
\qquad
\tau\ge 0,
\]
并定义达到水平 \(\rho\) 的最小归一化阈值
\[
\tau_\ast(\rho):=
\inf\left\{\tau\ge 0:\ \liminf_{m\to\infty}\Psi_m(\tau)\ge \rho\right\},
\qquad
0\le \rho\le 1.
\]
则有严格闭式
\[
\tau_\ast(\rho)=
\begin{cases}
\displaystyle \frac{\sqrt5}{\varphi^2}\,\rho,
& 0\le \rho\le \dfrac{\varphi}{\sqrt5},\\[3mm]
\displaystyle \frac{\sqrt5\,\rho-\varphi^{-1}}{\varphi},
& \dfrac{\varphi}{\sqrt5}\le \rho\le 1.
\end{cases}
\]
\end{theorem}
\begin{proof}
由定理 \ref{thm:xi-time-part70b-center-capacity-two-kink-skeleton}，
\[
\Psi_m(\tau)\longrightarrow \Psi(\tau):=
\frac{\varphi}{\sqrt5}\min\{1,\tau\}
+
\frac{1}{\sqrt5}\min\{\varphi^{-1},\tau\}.
\]
因此
\[
\tau_\ast(\rho)=\inf\{\tau\ge 0:\ \Psi(\tau)\ge \rho\}.
\]
而
\[
\Psi(\tau)=
\begin{cases}
\dfrac{\varphi^2}{\sqrt5}\tau, & 0\le \tau\le \varphi^{-1},\\[6pt]
\dfrac{\varphi}{\sqrt5}\tau+\dfrac{1}{\varphi\sqrt5}, & \varphi^{-1}\le \tau\le 1,\\[8pt]
1, & \tau\ge 1,
\end{cases}
\]
故广义逆由两段线性方程直接求得，即为所示公式。
\end{proof}

\begin{theorem}[escort 响应的逆温度阈值公式]\label{thm:xi-time-part70ab-escort-response-inverse-threshold-closed}
对固定 \(q\ge 0\)，定义 escort 归一化响应
\[
\Psi_{m,q}^{\mathrm{esc}}(\tau):=
\mathbb E_{\pi_{m,q}}
\left[
\min\!\left(1,\frac{\tau A_m}{d_m(W_m^{(q)})}\right)
\right],
\qquad
\tau\ge 0.
\]
再定义达到水平 \(\rho\) 的最小归一化阈值
\[
\tau_{q,\ast}(\rho):=
\inf\left\{\tau\ge 0:\ \liminf_{m\to\infty}\Psi_{m,q}^{\mathrm{esc}}(\tau)\ge \rho\right\},
\qquad
0\le \rho\le 1.
\]
记
\[
\rho_c(q):=\theta(q)+\varphi^{-1}\bigl(1-\theta(q)\bigr).
\]
则有闭式
\[
\tau_{q,\ast}(\rho)=
\begin{cases}
\displaystyle \frac{\rho}{1+(\varphi-1)\theta(q)},
& 0\le \rho\le \rho_c(q),\\[3mm]
\displaystyle \frac{\rho-\theta(q)}{1-\theta(q)},
& \rho_c(q)\le \rho\le 1.
\end{cases}
\]
\end{theorem}
\begin{proof}
由定理 \ref{thm:xi-time-part70b-escort-two-kink-response}，
\[
\Psi_{m,q}^{\mathrm{esc}}(\tau)\longrightarrow
\Psi_q^{\mathrm{esc}}(\tau):=
\bigl(1-\theta(q)\bigr)\min\{1,\tau\}
+
\theta(q)\min\{1,\varphi\tau\}.
\]
因此
\[
\tau_{q,\ast}(\rho)=\inf\{\tau\ge 0:\ \Psi_q^{\mathrm{esc}}(\tau)\ge \rho\}.
\]
把 \(\Psi_q^{\mathrm{esc}}\) 在 \([0,\varphi^{-1}]\) 与 \([\varphi^{-1},1]\) 上分段展开，得到
\[
\Psi_q^{\mathrm{esc}}(\tau)=
\begin{cases}
\bigl(1+(\varphi-1)\theta(q)\bigr)\tau,
& 0\le \tau\le \varphi^{-1},\\[6pt]
\theta(q)+\bigl(1-\theta(q)\bigr)\tau,
& \varphi^{-1}\le \tau\le 1,\\[6pt]
1, & \tau\ge 1.
\end{cases}
\]
其连接点函数值为
\[
\Psi_q^{\mathrm{esc}}(\varphi^{-1})
=
\theta(q)+\varphi^{-1}\bigl(1-\theta(q)\bigr)
=
\rho_c(q).
\]
故广义逆由上述两段线性函数分别求反即可。
\end{proof}

\begin{theorem}[escort 熵谱的完整常数项]\label{thm:xi-time-part70ab-escort-renyi-entropy-full-constant}
对固定 \(q\ge 0\)、\(a>0\) 且 \(a\neq 1\)，定义 escort Rényi 熵
\[
H_a(\pi_{m,q})
:=
\frac{1}{1-a}\log\sum_{w\in X_m}\pi_{m,q}(w)^a.
\]
则有渐近展开
\[
H_a(\pi_{m,q})
=
m\log\varphi
-\frac12\log 5
+
\frac{1}{1-a}
\log\!\Bigl(
\varphi^{1-a}\bigl(1-\theta(q)\bigr)^a+\theta(q)^a
\Bigr)
+o(1).
\]
当 \(a\to 1\) 时，得到 Shannon 熵常数项
\[
H(\pi_{m,q})
=
m\log\varphi
-\frac12\log 5
+
\bigl(1-\theta(q)\bigr)\log\varphi
+
h\bigl(\theta(q)\bigr)
+o(1),
\]
其中
\[
h(t):=-t\log t-(1-t)\log(1-t).
\]
\end{theorem}
\begin{proof}
记
\[
S_t(m):=\sum_{w\in X_m} d_m(w)^t.
\]
由定理 \ref{thm:xi-time-part70aa-binfold-all-real-moments-closed} 与 Binet 公式，对每个固定实数 \(t\) 都有
\[
S_t(m)
=
\frac{\varphi+\varphi^{-t}}{\sqrt5}
\bigl(2^t\varphi^{1-t}\bigr)^m
\bigl(1+o(1)\bigr).
\]
因此
\[
\sum_{w\in X_m}\pi_{m,q}(w)^a
=
\frac{S_{aq}(m)}{S_q(m)^a}
=
5^{(a-1)/2}\,
\varphi^{(1-a)m}\,
\frac{\varphi+\varphi^{-aq}}{(\varphi+\varphi^{-q})^a}
\bigl(1+o(1)\bigr).
\]
取对数并除以 \(1-a\)，得到
\[
H_a(\pi_{m,q})
=
m\log\varphi
-\frac12\log 5
+
\frac{1}{1-a}
\log\!\frac{\varphi+\varphi^{-aq}}{(\varphi+\varphi^{-q})^a}
+o(1).
\]
再由
\[
\theta(q)=\frac{\varphi^{-q}}{\varphi+\varphi^{-q}},
\qquad
1-\theta(q)=\frac{\varphi}{\varphi+\varphi^{-q}},
\]
可将最后一项重写为
\[
\frac{\varphi+\varphi^{-aq}}{(\varphi+\varphi^{-q})^a}
=
\varphi^{1-a}\bigl(1-\theta(q)\bigr)^a+\theta(q)^a.
\]
这就给出第一条公式。

当 \(a\to 1\) 时，对
\[
F_q(a):=
\log\!\Bigl(
\varphi^{1-a}\bigl(1-\theta(q)\bigr)^a+\theta(q)^a
\Bigr)
\]
应用 l'Hospital 法则，得到
\[
\lim_{a\to 1}\frac{F_q(a)}{1-a}
=
-F_q'(1)
=
\bigl(1-\theta(q)\bigr)\log\varphi+h\bigl(\theta(q)\bigr).
\]
代回即得 Shannon 常数项。
\end{proof}

\begin{theorem}[容量极限曲线的负矩 Mellin 反演]\label{thm:xi-time-part70ab-capacity-limit-negative-moment-mellin}
定义中心容量极限
\[
m_\ast(\tau):=
\lim_{m\to\infty}
\frac{\sqrt5}{\varphi^2}\,
2^{-m}\,\mathcal C_m^{\mathrm{bin,cont}}(\tau A_m),
\qquad
\tau\ge 0.
\]
则对任意实数 \(s>0\)，成立严格恒等式
\[
s(s+1)\int_0^\infty \tau^{-s-2}\bigl(\tau-m_\ast(\tau)\bigr)\,\dd\tau
=
\varphi^{-1}+\varphi^{s-2}.
\]
\end{theorem}
\begin{proof}
由定理 \ref{thm:xi-time-part70b-center-capacity-two-kink-skeleton}，
\[
m_\ast(\tau)=
\varphi^{-1}\min\{1,\tau\}
+
\varphi^{-2}\min\{\varphi^{-1},\tau\}.
\]
因而若定义双原子概率测度
\[
\mu_\ast:=\varphi^{-2}\delta_{\varphi^{-1}}+\varphi^{-1}\delta_1,
\]
则有
\[
m_\ast(\tau)=\int \min\{y,\tau\}\,\dd\mu_\ast(y).
\]
对任意固定 \(y>0\) 与 \(s>0\)，直接计算得
\[
\int_0^\infty \tau^{-s-2}\bigl(\tau-\min\{y,\tau\}\bigr)\,\dd\tau
=
\int_y^\infty \tau^{-s-2}(\tau-y)\,\dd\tau
=
\frac{y^{-s}}{s(s+1)}.
\]
两边乘以 \(s(s+1)\) 并对 \(\mu_\ast\) 积分，即得
\[
s(s+1)\int_0^\infty \tau^{-s-2}\bigl(\tau-m_\ast(\tau)\bigr)\,\dd\tau
=
\int y^{-s}\,\dd\mu_\ast(y)
=
\varphi^{-1}+\varphi^{-2}\varphi^s
=
\varphi^{-1}+\varphi^{s-2}.
\]
证毕。
\end{proof}

\begin{corollary}[Stirling--Bernoulli 几何系数的容量解释]\label{cor:xi-time-part70ab-bernoulli-hierarchy-capacity-interpretation}
\leanverified{paper\_xi\_time\_part70ab\_bernoulli\_hierarchy\_capacity\_interpretation}
对任意整数 \(r\ge 1\)，都有
\[
\varphi^{-1}+\varphi^{2r-3}
=
(2r-1)(2r)\int_0^\infty \tau^{-2r-1}\bigl(\tau-m_\ast(\tau)\bigr)\,\dd\tau.
\]
从而 Stirling--Bernoulli 层级中第 \(r\) 级系数可写为
\[
\frac{B_{2r}}{r(2r-1)}
\bigl(\varphi^{-1}+\varphi^{2r-3}\bigr)
=
\frac{2B_{2r}}{r}
\int_0^\infty \tau^{-2r-1}\bigl(\tau-m_\ast(\tau)\bigr)\,\dd\tau.
\]
特别地，对任意固定 \(R\ge 1\)，定理 \ref{thm:fold-bin-gauge-constant-stirling-bernoulli-hierarchy} 的展开可重写为
\[
\log C_m
=
\log(2\pi)
+
\sum_{r=1}^{R}
\frac{2B_{2r}}{r}
\left(
\int_0^\infty \tau^{-2r-1}\bigl(\tau-m_\ast(\tau)\bigr)\,\dd\tau
\right)
\left(\frac{\varphi}{2}\right)^{(2r-1)m}
+
O\!\left(\left(\frac{\varphi}{2}\right)^{(2R+1)m}\right).
\]
\end{corollary}
\begin{proof}
在定理 \ref{thm:xi-time-part70ab-capacity-limit-negative-moment-mellin} 中取 \(s=2r-1\) 即得第一式。再把该恒等式代入定理 \ref{thm:fold-bin-gauge-constant-stirling-bernoulli-hierarchy} 的系数表达，便得到后两式。
\end{proof}

\noindent
由此，末位 \(0/1\) 的二态骨架不只控制纤维多重度的最终排序与均匀/escort 采样下的分位相图，还把中心容量曲线的广义逆、escort 温度响应的广义逆以及 \(\log C_m\) 的 Stirling--Bernoulli 几何权重全部还原为同一组双原子数据。于是，阈值/采样轴、路径/相位轴与 Bernoulli 塔的解析层级便在同一条末位二态主线上闭合。

\endinput
